\documentclass[11pt,a4paper]{article}

\usepackage[utf8]{inputenc}
\usepackage[T1]{fontenc}
\usepackage[portuguese]{babel}
\usepackage{lmodern}
\usepackage{microtype}
\usepackage{amsmath,amssymb}
\usepackage{graphicx}
\usepackage[a4paper,margin=2.3cm]{geometry}
\usepackage{xcolor}
\usepackage{booktabs}
\usepackage{tabularx}
\usepackage{array}
\usepackage{enumitem}
\usepackage[most]{tcolorbox}
\usepackage[colorlinks=true,linkcolor=blue!55!black,urlcolor=blue!55!black]{hyperref}
\usepackage{titlesec}

\definecolor{azul}{HTML}{1E4E79}
\definecolor{verdee}{HTML}{16A34A}
\definecolor{laranja}{HTML}{B45309}

\titleformat{\section}{\Large\bfseries\color{azul}}{\thesection}{0.6em}{}
\titleformat{\subsection}{\large\bfseries\color{azul!85!black}}{\thesubsection}{0.6em}{}

\newtcolorbox{intuicao}{colback=verdee!6,colframe=verdee!60!black,
  title=\textbf{Intuição},fonttitle=\bfseries,boxrule=0.6pt,left=3mm,right=3mm,top=2mm,bottom=2mm}
\newtcolorbox{atencao}{colback=laranja!7,colframe=laranja!75!black,
  title=\textbf{Atenção / armadilha comum},fonttitle=\bfseries,boxrule=0.6pt,left=3mm,right=3mm,top=2mm,bottom=2mm}
\newtcolorbox{resumo}{colback=azul!6,colframe=azul!70!black,
  title=\textbf{Em uma frase},fonttitle=\bfseries,boxrule=0.6pt,left=3mm,right=3mm,top=2mm,bottom=2mm}

\newcommand{\fig}[2]{\begin{figure}[ht]\centering
  \includegraphics[width=\linewidth]{figures/#1}\caption{#2}\end{figure}}

\title{\textbf{Comprehend}\\[2mm]
\large Guia didático completo do projeto de Política Monetária:\\
um modelo Novo-Keynesiano mínimo calibrado para o Chile\\[3mm]
\normalsize Teoria, código e gráficos explicados do zero --- Aula 5 (Dynare e Calibração), páginas 26--55}
\author{Material de estudo}
\date{\today}

\begin{document}
\maketitle
\begin{abstract}
\noindent Este documento foi escrito para quem está começando. Ele explica, passo a passo e sem
pressupor conhecimento prévio, \emph{toda} a teoria por trás do projeto: o que é um modelo
Novo-Keynesiano (NK), o que cada equação significa, o que cada parâmetro faz, como o modelo é
resolvido no Dynare, como se leem todos os gráficos, como os dados do Chile entram e como se faz a
previsão das páginas 48--55. Sempre que aparece um termo técnico, ele é definido. As caixas verdes
trazem a intuição; as laranjas, as armadilhas; as azuis, o resumo de uma frase.
\end{abstract}

\tableofcontents
\bigskip

% ======================================================================
\section{Para que serve tudo isso}
% ======================================================================

\textbf{Política monetária} é o conjunto de decisões do Banco Central sobre a \textbf{taxa básica de
juros}. No Chile, essa taxa se chama \textbf{TPM} (\emph{Tasa de Política Monetaria}); no Brasil, é a
Selic. Subir o juro esfria a economia e tende a reduzir a inflação; baixar o juro aquece a economia.
O Banco Central mexe no juro tentando manter a \textbf{inflação} perto de uma \textbf{meta} (no Chile,
3\% ao ano) sem causar oscilações desnecessárias na \textbf{atividade econômica}.

Para estudar esse mecanismo de forma disciplinada, os economistas usam um \textbf{modelo}: um conjunto
pequeno de equações que descreve como juros, inflação e atividade se influenciam. O modelo deste
projeto é o \textbf{Novo-Keynesiano (NK) mínimo}, com apenas três equações e três variáveis. Ele é a
espinha dorsal de modelos muito maiores usados por bancos centrais de verdade.

\begin{intuicao}
Pense no modelo como um ``simulador de voo'' da economia. Ele não é a economia real, mas captura as
forças principais. Com ele podemos perguntar: \emph{``e se o Banco Central subir o juro?''} ou
\emph{``e se houver um choque de custos?''} --- e ver a resposta sem precisar fazer o experimento na
economia de verdade.
\end{intuicao}

\subsection{As três variáveis do modelo}
O modelo acompanha três grandezas, todas medidas trimestralmente:
\begin{itemize}[nosep]
  \item $x_t$ --- o \textbf{hiato do produto} (\emph{output gap}): a diferença percentual entre o PIB
        observado e o ``PIB potencial'' (o que a economia produziria sem pressões). $x_t>0$ significa
        economia ``aquecida''; $x_t<0$, economia ``abaixo do potencial''.
  \item $\pi_t$ --- a \textbf{inflação} (variação dos preços) no trimestre, medida como desvio da sua
        média/meta.
  \item $i_t$ --- a \textbf{taxa de juros nominal} de política (a TPM), também trimestral.
\end{itemize}

\begin{atencao}
``Hiato'' não é o nível do PIB, é o \emph{desvio} do PIB em relação à sua tendência. Um país rico e um
país pobre podem ter o mesmo hiato (zero) se ambos estiverem exatamente no seu potencial. O hiato mede
\emph{ciclo}, não \emph{tamanho}.
\end{atencao}

\subsection{Três ideias que aparecem o tempo todo}
Antes da matemática, três conceitos que voltam sempre:

\paragraph{1. Desvios em torno do equilíbrio (log-linearização).} O modelo não trabalha com os níveis
``brutos'' (PIB em reais, inflação em \%), mas com \textbf{desvios} em relação a um ponto de
equilíbrio chamado \textbf{estado estacionário}. Por isso as variáveis são pequenas e ficam ``em torno
de zero''. Trabalhar com desvios transforma equações curvas (não lineares) em equações de linha reta
(lineares), que são fáceis de resolver. É como usar uma régua para aproximar uma curva perto de um
ponto: funciona muito bem para pequenas oscilações.

\paragraph{2. Expectativas racionais (\emph{forward-looking}).} As pessoas e empresas tomam decisões
hoje pensando no \emph{futuro}. Uma empresa que fixa preço hoje pensa na inflação dos próximos
trimestres. No modelo, isso aparece como termos $x_{t+1}$ e $\pi_{t+1}$ (variáveis ``do futuro''). A
notação $\pi_{t+1}$ dentro de uma equação de hoje significa ``a inflação que se \emph{espera} para o
próximo trimestre''.

\paragraph{3. Choques.} São surpresas imprevisíveis que empurram a economia: um choque de demanda
($\varepsilon^x$), um choque de custos ($\varepsilon^\pi$) e um choque de política monetária
($\varepsilon^i$). Tudo o que é ``dinâmico'' no modelo vem da reação a esses choques.

\begin{resumo}
O modelo descreve como hiato, inflação e juros --- medidos como pequenos desvios do equilíbrio e
influenciados pelas expectativas de futuro --- reagem a três tipos de choque.
\end{resumo}

% ======================================================================
\section{O modelo, equação por equação}
% ======================================================================

O modelo NK mínimo tem três equações. Vamos dissecar cada uma. Em forma compacta:
\begin{align}
x_t      &= \mathbb{E}_t x_{t+1} - \tfrac{1}{\sigma}\big(i_t - \mathbb{E}_t\pi_{t+1} - r^\ast\big) + \varepsilon^x_t, \tag{IS}\\[1mm]
\pi_t    &= \beta\,\mathbb{E}_t\pi_{t+1} + \kappa\,x_t + \varepsilon^\pi_t, \tag{NKPC}\\[1mm]
i_t      &= \rho_i\,i_{t-1} + (1-\rho_i)\big(r^\ast + \phi_\pi\,\pi_t + \phi_x\,x_t\big) + \varepsilon^i_t. \tag{Taylor}
\end{align}
O símbolo $\mathbb{E}_t$ significa ``valor esperado, dado o que se sabe no trimestre $t$''.

\subsection{Equação IS --- o lado da demanda}
\[
x_t = \underbrace{\mathbb{E}_t x_{t+1}}_{\text{hiato futuro esperado}} - \tfrac{1}{\sigma}\big(\underbrace{i_t - \mathbb{E}_t\pi_{t+1}}_{\text{juro real esperado}} - \underbrace{r^\ast}_{\text{juro neutro}}\big) + \varepsilon^x_t.
\]
\textbf{O que ela diz:} a demanda hoje ($x_t$) depende de (a)~quanto se espera de demanda amanhã e
(b)~do \textbf{juro real} comparado ao \textbf{juro neutro}.

\begin{itemize}[nosep]
  \item \textbf{Juro real} $= i_t - \mathbb{E}_t\pi_{t+1}$: é o juro nominal \emph{descontada} a
        inflação esperada. Se o banco cobra 10\% e a inflação esperada é 6\%, o juro real é $\approx$
        4\%. O que importa para poupar/gastar é o juro real, não o nominal.
  \item $r^\ast$ (lê-se ``r-estrela''): o \textbf{juro neutro} --- o juro real que mantém a economia
        equilibrada (nem aquecendo, nem esfriando). Se o juro real está \emph{acima} de $r^\ast$, a
        política está apertada e a demanda cai ($x_t$ diminui).
  \item $\sigma$ (sigma): mede o quanto a demanda reage ao juro real. $1/\sigma$ é a sensibilidade.
        Com $\sigma$ grande, as pessoas quase não mudam o consumo quando o juro muda.
\end{itemize}

\begin{intuicao}
A equação IS é a regra do ``vale a pena gastar agora ou esperar?''. Juro real alto (acima do neutro)
$\Rightarrow$ vale a pena poupar $\Rightarrow$ consumo/demanda caem hoje $\Rightarrow$ $x_t$ cai. O
termo $\mathbb{E}_t x_{t+1}$ aparece porque a decisão de consumo é \emph{intertemporal}: comparo hoje
com o futuro.
\end{intuicao}

\subsection{Curva de Phillips Novo-Keynesiana (NKPC) --- o lado da oferta}
\[
\pi_t = \beta\,\mathbb{E}_t\pi_{t+1} + \kappa\,x_t + \varepsilon^\pi_t.
\]
\textbf{O que ela diz:} a inflação de hoje depende da inflação \emph{esperada} para amanhã e do
aquecimento da economia ($x_t$).

\begin{itemize}[nosep]
  \item $\beta$ (beta): o \textbf{fator de desconto}, um número pouco abaixo de 1 (aqui $\approx
        0{,}993$). Mede o peso que se dá ao futuro. Como as empresas são \emph{forward-looking}, a
        inflação esperada quase ``passa direto'' para a inflação de hoje.
  \item $\kappa$ (kappa): a \textbf{inclinação da curva de Phillips}. Mede o quanto o aquecimento
        ($x_t>0$) pressiona os preços. $\kappa$ grande $\Rightarrow$ economia aquecida vira inflação
        rapidamente; $\kappa$ pequeno $\Rightarrow$ os preços reagem pouco ao ciclo.
\end{itemize}

\paragraph{De onde vem $\kappa$? O modelo de Calvo.} Imagine que, a cada trimestre, só uma
\emph{fração} das empresas pode reajustar o preço (a probabilidade de \emph{não} reajustar é
$\theta$). Quem reajusta o faz pensando em todo o futuro. Desse mecanismo (chamado ``preços de
Calvo'') deduz-se uma fórmula heurística:
\[
\kappa \;\approx\; \frac{(1-\theta)(1-\beta\theta)}{\theta}\,\zeta,
\]
onde $\zeta$ agrega elasticidades e \emph{markups}. Quanto mais ``rígidos'' os preços (maior
$\theta$), menor o $\kappa$, e mais devagar a inflação responde. Na aula de calibração, tratamos
$\kappa$ como um parâmetro reduzido (escolhido diretamente), sem precisar abrir todo o Calvo.

\begin{atencao}
A NKPC \emph{não} tem um termo de inflação passada ($\pi_{t-1}$). Ela é puramente voltada para o
futuro. Consequência importante (que veremos nos gráficos): este modelo mínimo \textbf{não tem inércia
de inflação}. A inflação não ``gruda'' --- ela reverte rápido. Modelos maiores adicionam indexação
($\pi_{t-1}$) para corrigir isso.
\end{atencao}

\subsection{Regra de Taylor --- o comportamento do Banco Central}
\[
i_t = \rho_i\,i_{t-1} + (1-\rho_i)\big(r^\ast + \phi_\pi\,\pi_t + \phi_x\,x_t\big) + \varepsilon^i_t.
\]
\textbf{O que ela diz:} o Banco Central define o juro reagindo à inflação e ao hiato, mas com
\textbf{suavização} (não muda o juro de forma brusca).

\begin{itemize}[nosep]
  \item $\phi_\pi$ (phi-pi): quanto o juro sobe quando a inflação sobe. \textbf{Princípio de Taylor:}
        para a economia ser estável, $\phi_\pi$ precisa ser $>1$ (ver Seção~\ref{sec:det}). Aqui
        $\phi_\pi=1{,}75$.
  \item $\phi_x$ (phi-x): quanto o juro reage ao hiato. Mantido fixo em $0{,}5$.
  \item $\rho_i$ (rho-i): \textbf{suavização de juros} (\emph{interest-rate smoothing}). É a inércia: o
        juro de hoje ``puxa'' o de amanhã. $\rho_i=0{,}80$ a $0{,}93$ significa que o BC move o juro em
        passos pequenos e graduais, como de fato fazem os bancos centrais.
  \item $\varepsilon^i_t$: o \textbf{choque de política monetária} --- a parte do juro que \emph{não}
        é explicada pela regra (uma decisão ``surpresa'', mais apertada ou mais frouxa que o usual).
\end{itemize}

\begin{intuicao}
A regra de Taylor é o ``piloto automático'' do BC: se a inflação passa da meta, suba o juro mais que
proporcionalmente ($\phi_\pi>1$); se a economia esfria, alivie um pouco ($\phi_x>0$); e faça isso
gradualmente ($\rho_i$ alto), para não dar solavancos.
\end{intuicao}

\begin{resumo}
IS = demanda reage ao juro real; NKPC = inflação reage ao aquecimento e ao futuro; Taylor = o BC
move o juro contra a inflação e o hiato, devagar.
\end{resumo}

% ======================================================================
\section{Os parâmetros e a calibração}
% ======================================================================

\textbf{Calibrar} é escolher os números (os parâmetros) do modelo. A tarefa do trabalho (slide~26)
pede exatamente isto, para o Chile:

\begin{enumerate}[nosep]
  \item Fixar $r^\ast \in \{2\%,3\%,4\%\}$ ao ano e \textbf{derivar} $\beta$.
  \item Estimar $\rho_i$ por um \textbf{AR(1)} na taxa de política trimestral.
  \item Escolher $\kappa \in \{0{,}07,\,0{,}10,\,0{,}13\}$ e comparar as IRFs.
  \item Variar $\phi_\pi$ entre $1{,}3$ e $2{,}2$ (com $\phi_x=0{,}5$) e avaliar estabilidade.
  \item Calibrar os desvios-padrão dos choques ($\sigma_{\varepsilon x},\sigma_{\varepsilon \pi},
        \sigma_{\varepsilon i}$) para obter uma FEVD ``parecida com a de um relatório de BC''.
\end{enumerate}
Resultado esperado: uma tabela ``parâmetro $\leftrightarrow$ alvo'', as IRFs e os momentos.

\subsection{$r^\ast$ e $\beta$: o juro neutro e o desconto}
O fator de desconto $\beta$ e o juro neutro $r^\ast$ estão amarrados por
\[
\beta = \frac{1}{1+r^\ast_{\text{trim}}}, \qquad r^\ast_{\text{trim}} = (1+r^\ast_{\text{anual}})^{1/4}-1.
\]
\textbf{Por quê?} Em equilíbrio, poupar rende $r^\ast$; descontar o futuro por $\beta$ tem que
compensar exatamente esse rendimento. Daí $\beta(1+r^\ast)=1$. A conversão de anual para trimestral usa
raiz quarta porque há 4 trimestres no ano.

\begin{center}
\begin{tabular}{cccc}
\toprule
$r^\ast$ anual & $r^\ast$ trimestral & $\beta$ \\
\midrule
2\% & 0,496\% & 0,99506 \\
3\% & 0,742\% & 0,99264 \\
4\% & 0,985\% & 0,99024 \\
\bottomrule
\end{tabular}
\end{center}
O \emph{baseline} usa $r^\ast=3\%$, logo $\beta=0{,}99264$. (Arquivo: \texttt{rstar\_beta\_table.csv}.)

\subsection{O que cada parâmetro faz (resumo)}
\begin{center}
\begin{tabular}{lll}
\toprule
Símbolo & Nome & Papel \\
\midrule
$\sigma=1{,}0$ & elast.\ intertemporal inversa & sensibilidade da demanda ao juro real \\
$\beta=0{,}993$ & fator de desconto & peso do futuro (vem de $r^\ast$) \\
$\kappa=0{,}10$ & inclinação de Phillips & força do canal hiato $\to$ inflação \\
$\rho_i=0{,}93$ & suavização de juros & inércia da TPM \\
$\phi_\pi=1{,}75$ & resposta à inflação & ativismo do BC (princípio de Taylor) \\
$\phi_x=0{,}50$ & resposta ao hiato & reação do BC à atividade \\
\bottomrule
\end{tabular}
\end{center}

\subsection{Calibração \emph{versus} estimação}
Há duas formas de obter os números:
\begin{itemize}[nosep]
  \item \textbf{Calibração}: escolher valores ``razoáveis'' (da literatura, de relatórios, do bom
        senso). É o foco desta aula.
  \item \textbf{Estimação}: deixar os dados ``escolherem'' os valores por métodos estatísticos
        (econometria, Bayesiano). O projeto também faz isso, como \emph{checagem} (Seção~\ref{sec:econ}).
\end{itemize}
Seguindo a orientação de ``na dúvida, mostrar mais de uma versão'', o projeto roda os dois lados:
$\rho_i$ calibrado (0,80) e estimado (0,934); duas matrizes-alvo de FEVD; etc.

\subsection{Os choques e a FEVD-alvo}
Cada choque tem um tamanho típico (desvio-padrão). O \emph{baseline} usa
$\sigma_{\varepsilon x}=0{,}0050$, $\sigma_{\varepsilon \pi}=0{,}00277$,
$\sigma_{\varepsilon i}=0{,}00068$. Esses números foram \textbf{calibrados} para que a
\textbf{decomposição da variância} (FEVD, Seção~\ref{sec:fevd}) tenha um formato realista: a inflação
dominada por choques de custo, o juro pelo próprio choque, o hiato por demanda e política. Não são uma
decomposição histórica oficial do Chile --- são uma escolha de modelagem.

% ======================================================================
\section{Resolvendo e diagnosticando o modelo}
% ======================================================================

\subsection{Estado estacionário}
O \textbf{estado estacionário} é o ponto de repouso: sem choques, onde as variáveis ficam paradas.
Aqui, $x=0$, $\pi=0$ e $i=r^\ast$ (o juro fica no neutro). Em modelos lineares, o estado estacionário
serve apenas como ``ponto de apoio'': toda a ação acontece \emph{em torno} dele. Como os observáveis
do Chile estão em desvios da média, na prática trabalhamos com $x=\pi=i=0$.

\subsection{Expectativas racionais e a condição de Blanchard--Kahn}
\label{sec:det}
Modelos com $\mathbb{E}_t x_{t+1}$ e $\mathbb{E}_t\pi_{t+1}$ têm um desafio: as variáveis dependem do
futuro, que depende do futuro do futuro\ldots Existe solução única e estável? A resposta vem da
\textbf{condição de Blanchard--Kahn (BK)}.

A ideia: o sistema tem ``raízes'' (autovalores). Cada variável que ``salta'' (não tem passado, como
$\pi$ e $x$) precisa de exatamente uma raiz \emph{explosiva} (módulo $>1$) para ser ``amarrada''. Aqui
há \textbf{2 variáveis de salto} ($x,\pi$) e \textbf{1 variável de estado} ($i_{t-1}$). Logo,
precisamos de \emph{exatamente} 2 raízes explosivas e 1 estável. O Dynare confirma:
\[
\text{autovalores: } 0{,}656 \;(\text{estável}), \quad 1{,}04, \quad 1{,}379 \;(\text{explosivas}).
\]
\begin{itemize}[nosep]
  \item Raízes demais dentro do círculo $\Rightarrow$ \textbf{indeterminação} (infinitas soluções).
  \item Raízes de menos $\Rightarrow$ \textbf{explosão} (sem solução estável).
  \item O número certo $\Rightarrow$ solução \textbf{única e estável} (``determinação'').
\end{itemize}

\begin{intuicao}
O \textbf{princípio de Taylor} ($\phi_\pi>1$) é o que garante a determinação: se o BC reage à inflação
mais que proporcionalmente, as expectativas ficam ``ancoradas'' e o sistema converge. Se o BC reage
pouco ($\phi_\pi<1$), profecias auto-realizáveis de inflação tomam conta (indeterminação). Veremos
isso no mapa de determinação (Figura~\ref{fig:det}).
\end{intuicao}

\subsection{O que o Dynare imprime}
\textbf{Dynare} é o programa (rodando aqui no Octave) que resolve o modelo. Ao rodar, ele:
\begin{enumerate}[nosep]
  \item \emph{Pré-processa} o arquivo \texttt{.mod} (lê as equações, calcula derivadas).
  \item Calcula o \emph{estado estacionário} e verifica \emph{Blanchard--Kahn}.
  \item Imprime as \textbf{funções de política} (\emph{policy functions}): como cada variável reage à
        taxa passada e a cada choque.
  \item Reporta \textbf{momentos teóricos} (desvios-padrão, correlações), a \textbf{FEVD} e as
        \textbf{IRFs}.
\end{enumerate}
As funções de política são a ``solução'' do modelo: por exemplo, um choque de custo
($\varepsilon^\pi>0$) eleva a inflação ($+0{,}86$), eleva o juro ($+0{,}23$) e reduz o hiato
($-0{,}67$). Os arquivos ficam em \texttt{outputs/dynare/}\allowbreak
\texttt{baseline/}.

% ======================================================================
\section{Lendo os resultados: todos os gráficos}
% ======================================================================

\subsection{IRFs do baseline (os três choques)}
\textbf{IRF} (\emph{Impulse Response Function}, função de resposta ao impulso) responde: ``se eu der
\emph{um} choque hoje, o que acontece com cada variável nos próximos trimestres?''. É o gráfico mais
importante do modelo.

\fig{irf_baseline_all_shocks.png}{IRFs do baseline. Linhas: hiato, inflação, juro. Colunas: choque de
demanda, de custo, de política monetária.}

Lendo coluna por coluna:
\begin{itemize}[nosep]
  \item \textbf{Choque de demanda} ($\varepsilon^x>0$): o hiato sobe ($+0{,}4$~p.p.), a inflação sobe
        um pouco, e o BC reage subindo o juro; tudo volta ao normal em poucos trimestres.
  \item \textbf{Choque de custo} ($\varepsilon^\pi>0$): a inflação sobe ($+0{,}25$), mas o hiato
        \emph{cai} (o BC aperta para conter a inflação). Esse é o \textbf{trade-off} clássico: conter
        inflação custa atividade.
  \item \textbf{Choque de política} ($\varepsilon^i>0$, aperto surpresa): hiato \emph{e} inflação
        caem; o juro sobe e desce gradualmente (por causa da suavização $\rho_i$).
\end{itemize}

\begin{atencao}
As respostas convergem rápido, e algumas trocam de sinal por um trimestre (não são perfeitamente
monótonas). Isso é normal num modelo \emph{forward-looking} com suavização: há um pequeno
``\emph{undershooting}'' antes de voltar ao equilíbrio.
\end{atencao}

\subsection{Momentos e volatilidades: modelo \emph{versus} dados}
\textbf{Momentos} são estatísticas resumo: desvio-padrão (o quanto a variável oscila) e
autocorrelação (o quanto o valor de hoje se parece com o de ontem --- ``persistência'').

\fig{moments_model_vs_data.png}{Momentos do modelo \emph{versus} dados do Chile. Esquerda:
desvios-padrão. Direita: autocorrelação de 1ª ordem.}

Duas lições importantes:
\begin{itemize}[nosep]
  \item O modelo \textbf{subestima a volatilidade} dos dados (barras azuis/laranjas bem menores que as
        pretas/cinzas). A economia real teve choques enormes (crise de 2008, COVID-2020, inflação de
        2022) que um modelo mínimo não reproduz.
  \item A \textbf{persistência da inflação} nos dados ($\approx 0{,}4$) é muito maior que no modelo
        ($\approx 0{,}05$). É a confirmação numérica da armadilha que vimos: a NKPC pura não tem
        inércia. Esse é um limite conhecido do modelo mínimo.
\end{itemize}

\subsection{FEVD --- decomposição da variância}
\label{sec:fevd}
\textbf{FEVD} (\emph{Forecast Error Variance Decomposition}) responde: ``de toda a oscilação de uma
variável, quanto \% vem de cada choque?''. É como dividir a ``culpa'' pela variabilidade.

\fig{fevd_calibration_comparison.png}{FEVD: alvos declarados \emph{versus} alcançado pelo modelo, em
duas calibrações dos choques.}

No baseline (didático): o \textbf{hiato} vem de demanda ($\sim$45\%) e política ($\sim$45\%); a
\textbf{inflação} é dominada por choques de \textbf{custo} ($\sim$80\%); o \textbf{juro} pelo seu
próprio choque ($\sim$75\%). Esse padrão é coerente com a estrutura NK e ``parecido com um relatório
de BC''. A figura mostra que a calibração \emph{atinge} os alvos (barras ``Target'' e ``Achieved''
quase iguais).

\subsection{Sensibilidade a $\kappa$ (inclinação de Phillips)}
\fig{irf_kappa_comparison.png}{IRF da inflação a um choque, para $\kappa\in\{0{,}07,0{,}10,0{,}13\}$.}
Quanto maior $\kappa$, mais forte o repasse do aquecimento para a inflação: a curva de inflação
responde mais e o BC precisa reagir mais. O arquivo \texttt{irf\_kappa\_tradeoffs.png} resume o
\emph{trade-off} (impacto, persistência e custo em produto) para cada $\kappa$.

\fig{irf_kappa_tradeoffs.png}{Métricas de \emph{trade-off} para cada $\kappa$.}

\subsection{Sensibilidade a $\phi_\pi$ e o mapa de determinação}
\fig{irf_phi_pi_comparison.png}{IRFs variando $\phi_\pi$ entre 1,3 e 2,2 (com $\phi_x=0{,}5$).}
Um $\phi_\pi$ maior reduz o impacto inicial de um choque de custo sobre a inflação e acelera um pouco a
convergência, mas exige movimentos maiores de juros e produto.

\fig{irf_phi_pi_tradeoffs.png}{Métricas do \emph{trade-off} quando a reação à inflação
$\phi_\pi$ aumenta.}
O gráfico separa quatro ideias que não devem ser confundidas. Primeiro, uma regra mais agressiva
reduz o \textbf{pico inicial} da inflação. Segundo, isso exige maior \textbf{resposta do juro}.
Terceiro, o aperto produz maior deslocamento do \textbf{hiato do produto}, que é o custo real da
desinflação. Quarto, a inflação pode cruzar o zero e ficar temporariamente negativa
(\emph{undershooting}); por isso, a soma absoluta da resposta inflacionária não precisa cair
monotonicamente quando $\phi_\pi$ aumenta. ``Mais agressivo'' significa melhor contenção inicial,
mas não menor variabilidade em todas as métricas.

\fig{determinacy_map.png}{Mapa de determinação: número de raízes estáveis (barras) e raiz estável
dominante (linha) em função de $\phi_\pi$.}
\label{fig:det}
Este é, talvez, o gráfico mais bonito do projeto. À esquerda (vermelho, $\phi_\pi<1$): \textbf{2}
raízes estáveis --- raízes \emph{demais} $\Rightarrow$ \textbf{indeterminação}. À direita (verde,
$\phi_\pi\geq 1$): \textbf{1} raiz estável $\Rightarrow$ \textbf{determinação}. A linha (raiz dominante)
despenca de $\approx 1{,}0$ para $\approx 0{,}66$ ao cruzar $\phi_\pi=1$ e continua caindo: quanto mais
ativo o BC, mais rápida a convergência. É o \textbf{princípio de Taylor} desenhado.

\subsection{Persistência: $\rho_i$ estimado \emph{versus} calibrado}
\fig{irf_rho_comparison.png}{IRFs com $\rho_i$ estimado (0,934) e calibrado (0,80).}
Com $\rho_i$ mais alto, o juro se move mais devagar e os efeitos duram mais (respostas mais
``esticadas''). A suavização alta estimada para o Chile amplia a persistência dos efeitos monetários.

% ======================================================================
\section{Dados do Chile e econometria}
\label{sec:econ}
% ======================================================================

\subsection{Os dados}
Os dados são \textbf{oficiais do Banco Central do Chile (BCCh)}, trimestrais, de 2001Q1 a 2026Q1 (101
trimestres). Três séries:
\begin{itemize}[nosep]
  \item \textbf{TPM} (juro): média trimestral da taxa diária. Convertida para fração trimestral
        $i_q=(1+\text{TPM}/100)^{1/4}-1$.
  \item \textbf{IPC} (inflação): variações mensais compostas em inflação trimestral.
  \item \textbf{PIB}: o \textbf{hiato} é o ciclo do log-PIB pelo \textbf{filtro HP}
        (\emph{Hodrick--Prescott}, $\lambda=1600$), que separa tendência de ciclo.
\end{itemize}

\fig{data_overview.png}{Dados oficiais do Chile: TPM, inflação (IPC) e hiato do produto.}
Observe os episódios: inflação saltando em 2008 e em 2022 (até $\sim$15\% anualizado) e o tombo de
$\sim$15\% no hiato na COVID (2020). A linha pontilhada na inflação marca a meta de 3\%.

\paragraph{\emph{Demeaning} (tirar a média).} Para alimentar o modelo (cujo equilíbrio é zero), as
séries são transformadas em \textbf{desvios da própria média}: $\pi = \text{infl}_q-\overline{\text{infl}_q}$,
$i=i_q-\overline{i_q}$, $x=\text{hiato}-\overline{\text{hiato}}$. Importante: a média amostral da
inflação foi \textbf{3,79\%} ao ano (acima da meta de 3\%) e a da TPM, \textbf{4,09\%}. Por isso, no
modelo, ``inflação no equilíbrio'' corresponde a 3,79\%, e a meta de 3\% é um desvio \emph{negativo}.

\subsection{Checagens econométricas}
Além de calibrar, o projeto deixa os dados ``falarem'':
\begin{itemize}[nosep]
  \item \textbf{AR(1) da TPM}: estima a persistência $\rho_i$. Resultado: $\rho_i=0{,}934$ (erro-padrão
        HAC $0{,}049$). É um número alto --- a TPM é muito persistente. (\texttt{rhoi\_estimate.csv}.)
  \item \textbf{Regra de Taylor reduzida} (MQO e VI/2SLS): tenta recuperar $\phi_\pi,\phi_x$ dos dados.
        As estimativas são imprecisas (de 0,83 a 3,05 conforme o método) --- a evidência é ``ampla''.
  \item \textbf{Curva de Phillips}: estima $\kappa$ (em torno de 0,10 numa versão \emph{backward}).
  \item \textbf{Moda posterior bayesiana (DSGE)}: estima todos os parâmetros de uma vez por máxima
        verossimilhança penalizada por \emph{priors}. (\texttt{bayesian\_estimates.csv}.)
\end{itemize}

\begin{atencao}
O AR(1) mede ``persistência reduzida'', que mistura gradualismo de política, mudanças de regime e
variação do juro neutro. Ele \emph{não} isola o coeficiente estrutural de suavização. Por isso o
projeto reporta os dois valores ($0{,}80$ e $0{,}934$) e mostra os resultados sob ambos.
\end{atencao}

% ======================================================================
\section{Previsão: as páginas 48--55}
% ======================================================================

Esta é a parte que o trabalho enfatiza. Objetivo: usar o modelo \emph{calibrado} para
\textbf{projetar} $\pi$, $i$ e $x$ nos próximos 8 trimestres, a partir dos dados observados. A
ferramenta é o próprio \textbf{Dynare}.

\subsection{Como a previsão é feita (fielmente ao roteiro)}
O arquivo \texttt{dynare/nk\_chile\_forecast.mod} declara as variáveis observadas e roda:
\begin{center}\small
\texttt{varobs pi i x;}\\
\texttt{estimation(datafile='chile\_observables\_dynare.csv', mode\_compute=0, forecast=8);}
\end{center}
O que isso faz, em português:
\begin{itemize}[nosep]
  \item \texttt{varobs pi i x} --- diz ao Dynare que essas três séries são \emph{observadas}.
  \item \texttt{mode\_compute=0} --- \textbf{não estima nada}; mantém a calibração fixa.
  \item O Dynare roda o \textbf{filtro/suavizador de Kalman} sobre o histórico (uma técnica que ``lê''
        os choques que explicam os dados) e projeta 8 trimestres à frente.
  \item Devolve \texttt{oo\_.forecast.Mean} (a previsão central) e \texttt{HPDinf}/\texttt{HPDsup} (as
        bandas de incerteza).
\end{itemize}

\subsection{Previsão incondicional (o \emph{fan chart})}
``Incondicional'' = sem impor nada; deixar o modelo seguir suas próprias leis a partir do último
estado observado.

\fig{forecast_fanchart.png}{Previsão incondicional gerada pelo Dynare (\texttt{forecast=8}), com bandas
de credibilidade HPD, em níveis anualizados.}

Leitura: a partir do último trimestre, o modelo projeta a inflação indo de \textbf{3,30\%} (horizonte
1) a \textbf{3,76\%} (horizonte 8), convergindo para a média de 3,79\%; a TPM cai de 4,36\% para
$\approx$4,11\%; o hiato fecha em direção a zero. A \textbf{banda azul} mostra a incerteza: dado o
tamanho dos choques, é a faixa onde a variável provavelmente cai. Note o ``salto'' da inflação: ela
estava em 5,7\% e a projeção começa em 3,3\%.

\begin{atencao}
Por que esse salto? Porque o modelo \textbf{não tem inércia de inflação} (a NKPC é pura) e porque o
\textbf{único estado} que o modelo carrega é o juro passado $i_{t-1}$ --- ele ``esquece'' o último
hiato e a última inflação. O modelo trata os 5,7\% como um choque transitório e reverte rápido para a
média. É uma limitação real, não um erro de código. Validação: a média do Dynare bate, até o último
dígito, com a forma reduzida analítica (a raiz estável $a_i=0{,}656$).
\end{atencao}

\subsection{Previsões condicionais (cenários de política)}
``Condicional'' = \emph{impor} uma trajetória a alguma variável e perguntar o que acontece com as
outras. Usa o comando \texttt{conditional\_forecast} do Dynare. O slide propõe dois tipos: fixar o
\emph{juro} ou fixar a \emph{inflação}. Implementamos três cenários:

\fig{conditional_scenarios.png}{Três cenários condicionais (manter o juro, apertar, e forçar a
inflação à meta), ancorados no estado atual.}

\begin{enumerate}[nosep]
  \item \textbf{Manter a TPM em 4,5\% por 2 trimestres.} Como 4,5\% está acima da média (4,09\%), já é
        levemente contracionista: a inflação cai para \textbf{3,04\%} e o hiato para $-0{,}63\%$.
  \item \textbf{Apertar para $\approx$5,1\% por 4 trimestres.} Forte aperto: a inflação despenca para
        \textbf{1,91\%}, mas o hiato vai a $-1{,}60\%$. (A amplificação grande de fixar o juro num
        modelo \emph{forward-looking} é o ``\emph{forward guidance puzzle}''.)
  \item \textbf{Forçar a inflação à meta de 3\% por 4 trimestres.} Como a meta (3\%) está \emph{abaixo}
        da média (3,79\%), exige aperto: a TPM sobe a \textbf{4,52\%} e o hiato vai a $-0{,}67\%$.
\end{enumerate}

\paragraph{Quanto custa cada cenário? Os choques implícitos.} Para impor cada trajetória, o BC precisa
``empurrar'' o juro com choques $\varepsilon^i$. O gráfico abaixo mostra o tamanho desses empurrões:

\fig{conditional_policy_shocks.png}{Inovações de política $\varepsilon^i$ necessárias para cumprir cada
caminho condicional.}
O aperto a 5,1\% exige a maior inovação inicial; forçar a meta de 3\% exige um aperto pontual mais
modesto. Esses choques tornam \emph{explícito} o esforço de política embutido em cada cenário.

\begin{atencao}
Uma nuance técnica: o \texttt{conditional\_forecast} do Dynare parte do \emph{estado estacionário}
(onde $\pi$ já é a média de 3,79\%). Por isso o cenário de meta foi calibrado com o desvio correto
($3\% \Rightarrow \pi_{\text{dev}}=-0{,}00192$). Os cenários da figura, porém, são reancorados no
\emph{estado atual} pela forma reduzida exata (idêntica à solução do Dynare), para conectarem
continuamente com o histórico. Os dois métodos concordam no sinal e no mecanismo.
\end{atencao}

% ======================================================================
\section{Análise macro aprofundada: do modelo mínimo ao mundo real}
% ======================================================================

Esta seção responde a uma pergunta diferente: \emph{o que aprendemos quando confrontamos o modelo
com a história, com métodos econométricos alternativos e com canais que o baseline omite?} Cada
ferramenta tem um tipo de resposta e um limite próprio.

\subsection{Suavizador de Kalman e decomposição histórica}
O \textbf{filtro de Kalman} atualiza a estimativa do estado à medida que cada observação chega. O
\textbf{suavizador} usa a amostra inteira para revisar a estimativa de cada trimestre. Dado o modelo,
ele pergunta: quais sequências de $\varepsilon^x$, $\varepsilon^\pi$ e $\varepsilon^i$ seriam
necessárias para reproduzir $x$, $\pi$ e $i$?

\fig{history_shock_decomposition.png}{Decomposição histórica: contribuição de demanda, custo e
política para cada variável observada.}
\fig{history_smoothed_shocks.png}{Choques estruturais recuperados pelo suavizador de Kalman.}

\begin{atencao}
Temos três observáveis, três choques e nenhum erro de medida. Por isso, o modelo consegue reproduzir
os dados quase exatamente. Isso é uma ótima validação computacional, mas não é identificação causal:
se a inflação sobe, o modelo é obrigado a distribuir essa alta entre os poucos choques disponíveis.
\end{atencao}

\subsection{Contrafactual de política e crítica de Lucas}
Depois de recuperar os choques, alimentamos a mesma história em uma regra mais agressiva,
$\phi_\pi=2{,}5$. Isso responde: \emph{mantendo todo o resto igual, qual seria a trajetória?}

\fig{history_counterfactual.png}{Contrafactual com reação mais forte à inflação.}

O resultado é \textbf{condicional}, não uma prova do que ``teria acontecido''. Pela \textbf{crítica de
Lucas}, famílias, firmas, expectativas e até os choques mudam quando a regra muda. O contrafactual é
útil para entender mecanismos, não para julgar retrospectivamente o BCCh.

\subsection{SVAR, Cholesky e comparação com o DSGE}
Um \textbf{VAR} explica cada variável por defasagens de todas as variáveis. Para chamar uma inovação
de ``choque monetário'', precisamos de identificação. No \textbf{SVAR recursivo}, a decomposição de
\textbf{Cholesky} impõe a ordem $[x,\pi,i]$: dentro do trimestre, o juro pode reagir a atividade e
inflação, mas estas não reagem contemporaneamente ao choque de juros.

\fig{svar_vs_dsge_irf.png}{IRFs de um choque monetário no SVAR e no DSGE.}

O VAR(4) não rejeita resíduos brancos ($p=0{,}169$), mas rejeita normalidade e mostra um
\emph{activity puzzle}: a atividade sobe inicialmente após o aperto identificado. Isso pode refletir
antecipação do BC, variáveis omitidas, a ordem recursiva ou instabilidade da amostra. A comparação
não escolhe um vencedor; ela mostra quais restrições do DSGE não são reproduzidas automaticamente
pelos dados.

\subsection{$r^\ast$ variável no tempo e modelos de espaço de estados}
Em um \textbf{modelo de espaço de estados}, uma variável não observada evolui ao longo do tempo e
gera observações com ruído. Aqui, uma tendência local do juro real ex ante, controlada pelo hiato,
serve de proxy para $r^\ast$.

\fig{rstar_time_varying.png}{Proxy estatística de $r^\ast$ comparada à grade fixa de 2/3/4\%.}

A estimativa final é 0,85\% a.a., mas a série vai de $-4{,}72\%$ a 3,05\%. Esse intervalo não deve ser
tratado como trajetória estrutural verdadeira: inflação esperada por AR, prêmio de risco, filtro HP e
crises entram na medida. Ela ensina que $r^\ast$ é incerto e móvel.

\subsection{Previsão pseudo-fora-da-amostra}
Uma boa previsão deve enfrentar dados que não participaram da estimação. Usamos janela
\textbf{expansiva}: estimamos até uma data, prevemos 1 ou 4 trimestres e avançamos. Comparamos NK,
AR(1), passeio aleatório e VAR pelo \textbf{RMSE}, a raiz do erro quadrático médio.

\fig{forecast_oos_comparison.png}{RMSE do NK e dos benchmarks nos horizontes de 1 e 4 trimestres.}

O NK é pior no horizonte de um trimestre, mas supera o passeio aleatório para inflação e hiato em
quatro trimestres. A estrutura econômica ajuda no médio prazo; a rigidez do modelo atrapalha no
curto. Chamamos o teste de \emph{pseudo} porque os dados foram revisados e o hiato HP usa a amostra
completa.

\subsection{Economia pequena e aberta: câmbio, UIP e cobre}
Para o Chile, o câmbio e o cobre são canais centrais:
\begin{itemize}[nosep]
  \item \textbf{UIP} (paridade descoberta de juros): a diferença de juros se relaciona à depreciação
        esperada, sujeita a prêmio de risco.
  \item \textbf{\emph{Pass-through}}: uma depreciação encarece importados e pode elevar o IPC.
  \item \textbf{Cobre}: preço maior melhora renda externa e demanda, mas seu efeito fiscal e cambial
        pode operar em direções adicionais.
\end{itemize}

\fig{open_economy_data.png}{Câmbio CLP/USD e preço internacional do cobre usados na calibração.}
\fig{open_economy_irfs.png}{IRFs do modelo ampliado com câmbio e cobre.}
\fig{open_vs_closed_monetary_irf.png}{Transmissão monetária no modelo aberto e no fechado.}

Uma depreciação aumenta a inflação no impacto; cobre positivo expande atividade e pressiona juros.
O modelo satisfaz Blanchard--Kahn, mas seus coeficientes estruturais são ilustrativos. O objetivo é
mostrar por que o baseline fechado é incompleto, não alegar estimação definitiva do setor externo.

\subsection{NKPC híbrida e inércia da inflação}
A Phillips pura ``esquece'' $\pi_{t-1}$. A versão híbrida adiciona indexação:
\[
 \pi_t=\gamma_\pi\pi_{t-1}+(1-\gamma_\pi)
 \left(\beta E_t\pi_{t+1}+\kappa x_t\right)+\varepsilon^\pi_t .
\]
Com $\gamma_\pi=0{,}35$, parte das firmas reajusta olhando a inflação passada. Isso eleva a
autocorrelação modelada de aproximadamente 0,05 para 0,38.

\fig{hybrid_nkpc_irfs.png}{A Phillips híbrida prolonga a resposta da inflação.}

O ajuste melhora, mas $\gamma_\pi$ foi calibrado. Melhor ajuste estatístico não prova, sozinho, que
indexação seja o mecanismo causal correto.

\subsection{Bayes completo: MCMC, R-hat e ESS}
Bayes combina \textbf{prior} e \textbf{verossimilhança} para gerar a posterior. O
\textbf{Metropolis--Hastings} propõe valores, aceita ou rejeita cada passo e constrói uma cadeia cuja
distribuição limite é a posterior.
\begin{itemize}[nosep]
  \item \textbf{Burn-in}: sorteios iniciais descartados; aqui, 30\%.
  \item $\hat R$: compara variação entre e dentro das cadeias; próximo de 1 é desejável.
  \item \textbf{ESS}: tamanho efetivo após descontar autocorrelação; menor que o número bruto.
  \item \textbf{intervalo credível}: região que contém determinada massa posterior, condicional ao
        modelo, dados e priors.
\end{itemize}

\fig{mcmc_posterior_diagnostics.png}{Posteriores de duas cadeias de 10.000 propostas e diagnóstico
$\hat R$.}

Todos os $\hat R$ ficaram abaixo de 1,05, mas o ESS ficou entre cerca de 92 e 268. A posterior média
é $\kappa=0{,}095$, $\rho_i=0{,}873$, $\phi_\pi=1{,}168$ e $\phi_x=0{,}295$. O intervalo de 90\% de
$\phi_\pi$, $[0{,}986;1{,}453]$, encosta na região abaixo de um: a força da regra e a distância até a
fronteira de determinação continuam incertas.

\subsection{Comparação por verossimilhança marginal}
A posterior diz quais parâmetros são plausíveis \emph{dentro} de um modelo. Para comparar modelos,
precisamos da \textbf{verossimilhança marginal} ou \textbf{evidência}:
\[
 p(y\mid M)=\int p(y\mid\theta,M)\,p(\theta\mid M)\,d\theta .
\]
Ela calcula quão bem o modelo prevê os dados antes de conhecer os parâmetros, fazendo uma média sobre
o prior. Por isso penaliza modelos que só ajustam bem em uma região muito estreita: adicionar um
parâmetro não garante evidência maior.

Usamos os mesmos 101 trimestres, os mesmos observáveis $(\pi,i,x)$ e os mesmos priors comuns. A
\textbf{aproximação de Laplace} substitui a posterior perto do modo por uma normal e usa o Hessiano
$H$ para medir a curvatura:
\[
 \log p(y\mid M)\approx \log p(y,\hat\theta\mid M)
 +\frac{k}{2}\log(2\pi)-\frac12\log|H(\hat\theta)|.
\]

\fig{bayesian_model_comparison.png}{Evidência marginal relativa: NKPC prospectiva versus híbrida.}

A log evidência foi 1021,15 para a prospectiva e 1033,17 para a híbrida. A diferença de 12,02
log-pontos corresponde a um \textbf{fator de Bayes aproximado de 166 mil} a favor da híbrida, com
$\gamma_\pi=0{,}341$ no modo. Em palavras: mesmo pagando a penalização do parâmetro extra, a inflação
passada melhora muito o ajuste probabilístico.

\begin{atencao}
É evidência de Laplace, local e condicionada aos priors. Ela pode exagerar a precisão se a posterior
for assimétrica ou multimodal. Não rodamos uma segunda MCMC completa para a híbrida nem usamos
\emph{bridge sampling}; portanto, o número deve ser lido como comparação formal aproximada, não como
verdade absoluta.
\end{atencao}

% ======================================================================
\section{Limitações, glossário e mapa do código}
% ======================================================================

\subsection{Limitações (honestas)}
\begin{itemize}[nosep]
  \item O \textbf{baseline} é fechado; a extensão aberta ainda é pequena, calibrada e sem demanda
        externa, prêmio de risco estrutural ou bloco fiscal.
  \item \textbf{Único estado} ($i_{t-1}$): não carrega independentemente o último hiato e inflação.
  \item \textbf{Sem inércia de inflação}: reversão rápida demais (ver momentos).
  \item $r^\ast$ fixado em cenários; na realidade varia no tempo. Os choques são calibrados, não
        identificados historicamente.
  \item O \emph{outlier} da COVID distorce momentos e estimação.
  \item \textbf{Nada aqui é previsão oficial nem recomendação de política} --- é um laboratório
        estrutural pequeno e transparente.
\end{itemize}

\subsection{Glossário rápido de símbolos}
\begin{center}\small
\begin{tabular}{ll@{\qquad}ll}
\toprule
$x_t$ & hiato do produto & $\rho_i$ & suavização de juros \\
$\pi_t$ & inflação & $\phi_\pi,\phi_x$ & respostas da regra de Taylor \\
$i_t$ & juro nominal (TPM) & $r^\ast$ & juro real neutro \\
$\beta$ & fator de desconto & $\varepsilon^x,\varepsilon^\pi,\varepsilon^i$ & choques \\
$\sigma$ & elast.\ intertemporal inv. & $\mathbb{E}_t$ & expectativa em $t$ \\
$\kappa$ & inclinação de Phillips & IRF/FEVD & resposta a impulso / decomp.\ de variância \\
\bottomrule
\end{tabular}
\end{center}

\subsection{Mapa do código (o que cada arquivo faz)}
\begin{center}\scriptsize
\begin{tabularx}{\textwidth}{@{}>{\raggedright\arraybackslash}p{0.39\textwidth}X@{}}
\toprule
Arquivo & Função \\
\midrule
\texttt{build\_chile\_dataset.py} & baixa/monta os dados do BCCh e os observáveis \\
\texttt{estimate\_rhoi\_chile.py} & AR(1) da TPM ($\rho_i$) \\
\texttt{estimate\_taylor\_rule.py} & regra de Taylor por MQO e VI \\
\texttt{estimate\_nkpc.py} & inclinação $\kappa$ \\
\texttt{calibrate\_shocks.py} & desvios-padrão dos choques para a FEVD-alvo \\
\texttt{generate\_dynare\_models.py} & gera o \emph{baseline} e os 18 cenários \texttt{.mod} \\
\texttt{run\_dynare\_batch.py} & roda os \texttt{.mod} no Octave/Dynare \\
\texttt{run\_bayesian.py} & estimação por moda posterior (DSGE) \\
\texttt{run\_forecast.py} & \textbf{previsão nativa do Dynare} (\texttt{forecast=8} + condicional) \\
\texttt{forecast\_model.py} & \emph{fan chart} em nível + cenários ancorados + choques \\
\texttt{run\_history.py / plot\_history.py} & decomposição histórica e contrafactual \\
\texttt{analyze\_svar.py} & SVAR recursivo e comparação com DSGE \\
\texttt{estimate\_time\_varying\_rstar.py} & proxy de $r^\ast$ em espaço de estados \\
\texttt{evaluate\_forecasts.py} & avaliação pseudo-fora-da-amostra \\
\texttt{run\_macro\_extensions.py} & modelos aberto e Phillips híbrida no Dynare \\
\texttt{run\_mcmc.py / plot\_mcmc.py} & posterior completa e diagnósticos das cadeias \\
\texttt{compare\_bayesian\_models.py} & evidência marginal de Laplace entre as duas NKPCs \\
\texttt{plot\_irfs.py} & figuras de IRFs e sensibilidade \\
\texttt{build\_final\_html.py} & monta o relatório \texttt{Entrega Final.html} \\
\texttt{nk\_chile\_forecast.mod} & modelo de previsão (\texttt{varobs}, \texttt{estimation}, \texttt{conditional\_forecast}) \\
\bottomrule
\end{tabularx}
\end{center}

\begin{resumo}
O projeto calibra um modelo NK mínimo para o Chile, resolve-o no Dynare, verifica determinação,
interpreta IRFs/momentos/FEVD, testa sensibilidades, confronta com econometria e, por fim, faz a
previsão nativa do Dynare das páginas 48--55. As extensões aprofundam história macro, validação,
economia aberta, inércia e incerteza bayesiana. Tudo é reprodutível e honesto quanto aos limites.
\end{resumo}

\end{document}
